% Standalone abbreviation and acronym entries moved from the alphabetical dictionary.
% Each definition provides the full form and, where necessary, context-dependent meanings.
\small
\begin{description}
\item[\textbf{A.C.}] A prescription abbreviation meaning “before meals,” from the Latin \textit{ante cibum}.

\item[\textbf{ALA}] An abbreviation for alpha-linolenic acid, an essential omega-3 fatty acid, or aminolevulinic acid, a heme-pathway compound used topically in some photodynamic treatments; context determines the meaning.

\item[\textbf{AP}] An abbreviation commonly meaning anteroposterior, or directed from front to back; in some contexts it may have another medical expansion.

\item[\textbf{ATM}] An abbreviation that may mean the temporomandibular joint or ataxia-telangiectasia mutated, a protein kinase involved in DNA-damage repair; context determines the meaning.

\item[\textbf{BF}] Physician's shorthand for Black female; demographic abbreviations should be used cautiously and replaced with precise, patient-preferred language when possible.

\item[\textbf{BPD}] An abbreviation commonly meaning bronchopulmonary dysplasia, chronic lung disease after premature birth; it can also mean borderline personality disorder in psychiatric contexts.

\item[\textbf{CA}] A medical abbreviation that may mean cancer, carcinoma, or calcium; the intended meaning depends on the surrounding clinical context.

\item[\textbf{CKD}] Abbreviation for Chronic Kidney Disease: a progressive loss of kidney function over
months to years, defined by a glomerular filtration rate (GFR) less than 60
$\mathrm{mL/min/1.73m^2}$ or markers of kidney damage (albuminuria, haematuria, structural
abnormalities) persisting for at least 3 months.

\item[\textbf{CLA}] An abbreviation that commonly means conjugated linoleic acid, a group of fatty-acid isomers; in other medical contexts it may have a different expansion.

\item[\textbf{CTG}] Abbreviation for cardiotocography; continuous electronic monitoring of fetal heart rate and uterine contractions during labour.

\item[\textbf{CVA}] An abbreviation for cerebrovascular accident, meaning a stroke caused by interruption of blood flow to the brain or
bleeding within the brain. It may cause sudden neurologic deficits.

\item[\textbf{CVS}] An abbreviation that may mean cardiovascular system or chorionic-villus sampling, a prenatal diagnostic procedure; context determines the meaning.

\item[\textbf{DDAVP}] The abbreviation for desmopressin, a synthetic vasopressin-like medicine used for some bleeding disorders, central diabetes insipidus, and nighttime urination.

\item[\textbf{DES}] An abbreviation that may mean diethylstilbestrol, a synthetic estrogen, or drug-eluting stent; diethylstilbestrol exposure during pregnancy is associated with later reproductive and cancer risks.

\item[\textbf{DHHS}] The abbreviation for the U.S. Department of Health and Human Services.

\item[\textbf{DNR}] An abbreviation for do-not-resuscitate, indicating that cardiopulmonary resuscitation should not be attempted if breathing or heartbeat stops.

\item[\textbf{DNS}] An abbreviation that may mean deviated nasal septum or dextrose in normal saline; the intended expansion depends on the clinical context.

\item[\textbf{DPT}] An older abbreviation for the combined diphtheria, pertussis, and tetanus vaccine; current formulations include DTaP and Tdap.

\item[\textbf{DT}] An abbreviation that may mean delirium tremens or diphtheria--tetanus vaccine; the intended meaning depends on context.

\item[\textbf{DTC}] An abbreviation that may mean differentiated thyroid carcinoma or direct-to-consumer, such as direct-to-consumer genetic testing; the intended meaning depends on context.

\item[\textbf{DTP}] An older abbreviation for diphtheria, tetanus, and pertussis vaccine; current childhood formulations are generally called DTaP.

\item[\textbf{DV}] An abbreviation that may mean domestic violence, ductus venosus, or another context-specific term; the full expression should be specified.

\item[\textbf{DYS}] An abbreviation or word root related to dysplasia, dystonia, dyslexia, or another condition beginning with dys-; the intended meaning depends on context.

\item[\textbf{EC}] An abbreviation that may mean enteric-coated, emergency contraception, or endothelial cell; the intended meaning depends on context.

\item[\textbf{EFT}] An abbreviation whose medical meaning depends on context, such as endocrine-function testing or emotional freedom technique; the intended expansion should be specified.

\item[\textbf{ENT}] Abbreviation for ear, nose, and throat; it commonly refers to the otolaryngology specialty or an otolaryngologist.

\item[\textbf{EOS}] Commonly an abbreviation for eosinophils or an eosinophil count in laboratory reports; the intended meaning should be confirmed from the report's context.

\item[\textbf{EPA}] Abbreviation for eicosapentaenoic acid, an omega-3 polyunsaturated fatty acid found in oily fish and fish-oil supplements.

\item[\textbf{EPR}] Abbreviation for electrophrenic respiration, also called diaphragm pacing, in which electrical stimulation helps activate the diaphragm in selected patients with respiratory muscle paralysis.

\item[\textbf{EPT}] Abbreviation for estrogen/progestin therapy, treatment using both an estrogen and a progestin, commonly discussed in menopausal hormone therapy when the uterus is present.

\item[\textbf{ER}] Abbreviation for emergency room or emergency department, the hospital service that evaluates and treats urgent illness and injury.

\item[\textbf{FCC}] In medicine, an abbreviation for familial colorectal cancer; in other contexts it may have nonmedical meanings.

\item[\textbf{FFI}] Abbreviation for fatal familial insomnia, a rare inherited prion disease causing progressively severe insomnia, autonomic dysfunction, and neurologic decline.

\item[\textbf{FISH}] Abbreviation for fluorescence in situ hybridization, a test that uses fluorescent DNA probes to detect specific genes, chromosomes, or chromosomal rearrangements in cells.

\item[\textbf{FLAP}] Abbreviation for 5-lipoxygenase-activating protein, a protein encoded by \textit{ALOX5AP} that helps produce leukotrienes involved in inflammation.

\item[\textbf{GALT}] Abbreviation for gut-associated lymphoid tissue, immune tissue in the gastrointestinal tract that samples intestinal contents and helps defend against pathogens.

\item[\textbf{GBV-C}] Abbreviation for GB virus C, also called human pegivirus, a blood-borne virus that can persist in people but is not established as a cause of chronic hepatitis.

\item[\textbf{GD}] In the RxList medical-dictionary context, abbreviation for Gaucher disease, an inherited lysosomal storage disorder caused by deficient glucocerebrosidase.

\item[\textbf{GI}] Abbreviation for gastrointestinal, relating to the digestive tract from the mouth through the anus and its associated organs.

\item[\textbf{GPO}] Abbreviation for group purchasing organization, an organization that combines buyers' purchasing power to negotiate prices and contracts for medicines, supplies, or services.

\item[\textbf{HAART}] Abbreviation for highly active antiretroviral therapy, an older term for combination antiretroviral treatment used to suppress HIV replication and preserve immune function.

\item[\textbf{HAPE}] Abbreviation for high-altitude pulmonary edema, a life-threatening accumulation of fluid in the lungs after rapid ascent; immediate descent and supplemental oxygen are urgent treatments.

\item[\textbf{HAV}] Abbreviation for hepatitis A virus, an RNA virus spread mainly by the fecal--oral route that causes acute hepatitis but usually does not cause chronic infection.

\item[\textbf{HBV}] Abbreviation for hepatitis B virus, a DNA virus transmitted through infected blood or body fluids that can cause acute or chronic hepatitis and cirrhosis or liver cancer.

\item[\textbf{HEPA}] Abbreviation for high-efficiency particulate air, describing a filter designed to remove a very high proportion of airborne particles of specified sizes from an air stream.

\item[\textbf{HGD}] Commonly an abbreviation for high-grade dysplasia, a marked abnormality of cell development that may be a precursor to invasive cancer but is not itself invasive cancer.

\item[\textbf{HNP}] Abbreviation for herniated nucleus pulposus, in which the inner material of an intervertebral disk protrudes through the outer ring and may compress a nerve root.

\item[\textbf{HPS}] Common medical abbreviations include hantavirus pulmonary syndrome, a severe respiratory illness, and Hermansky--Pudlak syndrome, an inherited disorder with oculocutaneous albinism and platelet-storage abnormalities; context determines the meaning.

\item[\textbf{HS}] Commonly an abbreviation for hereditary spherocytosis, an inherited hemolytic anemia; it can also mean half-strength, heart sounds, hamstrings, heparan sulfate, hernia sac, or house surgeon, so context is essential.

\item[\textbf{HSA}] Abbreviation that may mean a US health savings account or human serum albumin, the main plasma protein that helps maintain oncotic pressure; context determines the meaning.

\item[\textbf{ICU}] Abbreviation for intensive care unit, a hospital unit providing continuous monitoring and specialized support for critically ill patients.

\item[\textbf{ID}] A context-dependent abbreviation that may mean infectious disease, infectious diseases specialist, intellectual disability, or identification; the surrounding clinical text determines the meaning.

\item[\textbf{IFA}] Abbreviation for indirect fluorescent antibody, a laboratory method that detects antibodies or antigens using fluorescent labels, commonly in serologic testing.

\item[\textbf{IHS}] Abbreviation for the Indian Health Service, a US federal agency providing health services to American Indians and Alaska Native people and communities.

\item[\textbf{IPT}] Abbreviation for interpersonal therapy, a structured psychotherapy that focuses on relationships, grief, role transitions, and interpersonal difficulties.

\item[\textbf{IPV}] Abbreviation for inactivated polio vaccine, an injectable vaccine containing killed poliovirus and used to prevent poliomyelitis.

\item[\textbf{IQ}] Abbreviation for intelligence quotient, a standardized score from cognitive testing; it describes performance on specified tests and is not a complete measure of a person's abilities or worth.

\item[\textbf{ITC}] Abbreviation for in-the-canal, describing a hearing aid or other device positioned within the ear canal.

\item[\textbf{ITN}] Abbreviation for insecticide-treated bednet, a net treated with insecticide to reduce mosquito bites and malaria transmission.

\item[\textbf{IVC}] Commonly abbreviation for inferior vena cava, the large vein returning blood from the lower body to the heart; it can also mean intravenous cholangiogram, an X-ray study of the biliary system.

\item[\textbf{JIP}] Abbreviation for juvenile intestinal polyposis, an inherited condition in which children develop multiple hamartomatous polyps in the gastrointestinal tract and may have increased cancer risk.

\item[\textbf{KB}] In medical usage, an abbreviation for keratoderma blennorrhagicum or kilobase of DNA; lower-case \textit{kb} is preferred for kilobase, so context is important.

\item[\textbf{KUB}] Abbreviation for kidney, ureter, and bladder, commonly referring to a plain abdominal X-ray used to assess urinary tract calculi, bowel gas patterns, or other abdominal findings.

\item[\textbf{L}] Abbreviation for liter, a metric unit of volume equal to 1,000 milliliters.

\item[\textbf{LAVH}] Abbreviation for laparoscopically assisted vaginal hysterectomy, removal of the uterus through the vagina with laparoscopic assistance.

\item[\textbf{LL}] Commonly an abbreviation for lower lid of the eye or lower limit; the intended meaning depends on the clinical context.

\item[\textbf{LLL}] Abbreviation for the left lower lobe of the lung.

\item[\textbf{LLQ}] Abbreviation for left lower quadrant, usually referring to the lower-left quarter of the abdomen.

\item[\textbf{LMP}] Abbreviation for last menstrual period, the first day of the most recent menstrual bleeding, used in estimating gestational age when the date is reliable.

\item[\textbf{LPA}] LPA abbreviates Little People of America and may also refer to the left pulmonary artery.

\item[\textbf{LUL}] Abbreviation for the left upper lobe of the lung.

\item[\textbf{LUQ}] Abbreviation for the left upper quadrant of the abdomen.

\item[\textbf{LVF}] Abbreviation for left ventricular failure, meaning failure of the left side of the heart; it may also mean left ventricular function, depending on context.

\item[\textbf{MAOI}] Abbreviation for monoamine oxidase inhibitor, a drug that inhibits monoamine oxidase and is used in selected treatment of depression and other disorders.

\item[\textbf{MAP}] Common abbreviation for \textit{Mycobacterium avium} subspecies \textit{paratuberculosis}, an organism of the \textit{M. avium} complex; the abbreviation can have other meanings in different clinical contexts.

\item[\textbf{MCD}] An abbreviation with several medical meanings, including minimal change disease, mean corpuscular diameter, and medullary cystic disease; the intended meaning depends on context.

\item[\textbf{MCP}] Membrane cofactor protein, also called CD46, is a cell-surface protein that regulates complement activation and helps protect host cells from complement-mediated injury; the abbreviation can have other meanings in different contexts.

\item[\textbf{MDI}] An abbreviation that may mean metered-dose inhaler, multiple daily injections, or mental developmental index; the intended meaning depends on context.

\item[\textbf{MEDVAC}] Abbreviation for medical evacuation: the transport of ill or injured people to a medical facility or to a location providing a higher level of care.

\item[\textbf{MI}] Abbreviation for myocardial infarction, or heart attack: death of heart-muscle tissue caused by prolonged inadequate blood flow, usually from acute coronary artery blockage.

\item[\textbf{MM}] An abbreviation that may mean meningomyelocele, malignant melanoma, or multiple myeloma; in lowercase, \textit{mm} commonly means millimeter. The intended meaning depends on context.

\item[\textbf{MP}] Abbreviation for mercaptopurine, a purine-antimetabolite drug, or for molecular pathology; the intended meaning depends on context.

\item[\textbf{MPS}] Abbreviation for mucopolysaccharidosis, a group of inherited lysosomal storage disorders caused by deficient breakdown of glycosaminoglycans.

\item[\textbf{MSG}] Abbreviation for monosodium glutamate, the sodium salt of glutamic acid used as a flavor enhancer and contributing umami taste.

\item[\textbf{MTX}] Abbreviation for methotrexate, a folate-antagonist drug used in selected cancers, autoimmune diseases, and other conditions.

\item[\textbf{MVP}] Abbreviation for mitral valve prolapse, in which one or both mitral-valve leaflets bow into the left atrium during ventricular systole; it is often mild but can sometimes cause mitral regurgitation or arrhythmias.

\item[\textbf{MYH}] MYH is a gene-family abbreviation involving myosin or DNA-repair proteins; the specific gene and variant determine the clinical meaning.

\item[\textbf{ND}] ND is a context-dependent abbreviation that may mean naturopathic doctor, neurodevelopmental, or another term; the full meaning should be specified.

\item[\textbf{NEP}] NEP is a context-dependent abbreviation that may mean neuroendocrine pancreatic tumor, nasoendotracheal procedure, or another clinical term.

\item[\textbf{NF}] NF is a context-dependent abbreviation, including neurofibromatosis, a group of genetic disorders causing nerve-sheath tumors and other features.

\item[\textbf{NG}] NG commonly means nasogastric, referring to a tube passed through the nose into the stomach, but the abbreviation has other uses.

\item[\textbf{NS}] NS is a context-dependent abbreviation, commonly normal saline, a sterile sodium-chloride solution, or not significant; the full term should be stated.

\item[\textbf{OB}] OB is an abbreviation commonly meaning obstetrics, the care of pregnancy, labor, birth, and the postpartum period, but it may have other meanings.

\item[\textbf{OBS}] OBS is a context-dependent abbreviation that may mean observation, obstetric service, or organic brain syndrome; the intended term should be written out.

\item[\textbf{OCG}] OCG is a context-dependent abbreviation that may refer to oral contraceptive guidance, oculocardiac testing, or another clinical term; the full term should be specified.

\item[\textbf{OMS}] An abbreviation for oral and maxillofacial surgery; the intended meaning depends on the clinical context.

\item[\textbf{PAP}] PAP is an abbreviation with context-dependent meanings, including Papanicolaou testing, pulmonary alveolar proteinosis, or prostatic acid phosphatase; the full term should be specified.

\item[\textbf{PAS}] PAS is an abbreviation that may mean periodic-acid–Schiff staining, para-aminosalicylic acid, or pulmonary-artery systolic pressure; context determines the meaning.

\item[\textbf{PAT}] PAT is a context-dependent abbreviation that may refer to paroxysmal atrial tachycardia, a rapid rhythm beginning in the atria, or another clinical term.

\item[\textbf{PCA}] PCA is a context-dependent abbreviation, commonly meaning patient-controlled analgesia, in which a patient activates prescribed doses from an infusion device.

\item[\textbf{PCT}] PCT is a context-dependent abbreviation that may mean procalcitonin, porphyria cutanea tarda, or patient-care technician; the full term should be specified.

\item[\textbf{PDR}] PDR commonly means Physicians’ Desk Reference, a drug-information reference, but the abbreviation has other uses and should be interpreted in context.

\item[\textbf{PEP}] PEP is a context-dependent abbreviation, including post-exposure prophylaxis, a preventive medicine course after a possible infection exposure.

\item[\textbf{PHS}] PHS is an abbreviation whose meaning depends on context, including public health service or primary hyperhidrosis; the full term should be specified.

\item[\textbf{PI}] PI is a context-dependent abbreviation, including principal investigator, pulmonary insufficiency, or protease inhibitor; the intended meaning must be specified.

\item[\textbf{PLP}] PLP commonly means pyridoxal 5$^{\prime}$-phosphate, the active coenzyme form of vitamin B6 required by many enzymes, but the abbreviation has other uses.

\item[\textbf{PMD}] PMD is a context-dependent abbreviation, including Pelizaeus–Merzbacher disease, a rare inherited disorder of central myelination, or a professional medical designation.

\item[\textbf{PON}] PON is a context-dependent abbreviation that may refer to postoperative nausea, paraoxonase, or another clinical term; the full meaning should be stated.

\item[\textbf{PTU}] PTU is the abbreviation for propylthiouracil, an antithyroid medicine used selectively because of liver toxicity and other safety considerations.

\item[\textbf{PV}] PV is a context-dependent abbreviation, including polycythemia vera, plasma volume, or peripheral vision; the full term should be specified.

\item[\textbf{PVB}] PVB is a context-dependent abbreviation that may mean paravertebral block, a regional anesthetic technique, or another medical term.

\item[\textbf{PX}] PX is a context-dependent abbreviation that may mean prognosis, patient, prescription, or a genetic marker; clinical documentation should spell out the intended term.

\item[\textbf{Q.2H}] A prescription abbreviation meaning “every 2 hours.” Medication instructions should use clear wording because abbreviation misinterpretation can cause dosing errors.

\item[\textbf{Q.3H.}] A prescription abbreviation meaning “every 3 hours.” Medication instructions should use clear wording because abbreviation misinterpretation can cause dosing errors.

\item[\textbf{Q.S.}] A prescription abbreviation from the Latin “quantum sufficiat,” meaning “quantity sufficient” to make the required amount or concentration. The exact quantity must be specified by the formulation or prescription context.

\item[\textbf{R/O}] R/O is an abbreviation for rule out, indicating that a condition is being considered but has not been confirmed; clear documentation should state the actual diagnostic question.

\item[\textbf{RAD}] Acronym whose medical meaning depends on context, commonly radiation absorbed dose or right axis deviation on an ECG.

\item[\textbf{RBC}] RBC is the abbreviation for red blood cell, the cell that transports oxygen with hemoglobin and carries carbon dioxide back toward the lungs.

\item[\textbf{RF}] RF is a context-dependent abbreviation, including rheumatoid factor, radiofrequency, or risk factor; the complete term should be stated.

\item[\textbf{RLL}] RLL is the abbreviation for right lower lobe, usually referring to the lower lobe of the right lung.

\item[\textbf{RLQ}] RLQ is the abbreviation for right lower quadrant of the abdomen, containing structures including the appendix, cecum, and right adnexa.

\item[\textbf{RML}] RML is the abbreviation for right middle lobe, the middle lobe of the right lung.

\item[\textbf{RP}] RP is a context-dependent abbreviation, commonly retinitis pigmentosa, an inherited retinal degeneration causing night blindness and progressive visual-field loss.

\item[\textbf{RUQ}] RUQ is the abbreviation for right upper quadrant of the abdomen, containing the liver, gallbladder, duodenum, and portions of bowel and kidney.

\item[\textbf{SGOT}] SGOT is the former abbreviation for AST, an enzyme measured in blood that can rise with injury to liver, heart, muscle, or other tissues.

\item[\textbf{SGPT}] SGPT is the former abbreviation for ALT, a liver-associated enzyme that rises in many forms of hepatocellular injury.

\item[\textbf{SHP}] SHP is a context-dependent abbreviation that may refer to a healthcare professional, a gene or protein, or another clinical term; the full meaning should be specified.

\item[\textbf{SLV}] SLV is a context-dependent abbreviation that may mean single-lung ventilation, serum lithium value, or another term; the full meaning should be specified.

\item[\textbf{STD(S)}] Abbreviation for sexually transmitted disease(s), now often called sexually transmitted infections.

\item[\textbf{STS}] STS is a context-dependent abbreviation that may mean soft-tissue sarcoma, sodium thiosulfate, or another clinical term; the full term should be stated.

\item[\textbf{TMR}] TMR is a context-dependent abbreviation, including targeted muscle reinnervation, a surgical technique used to improve control of a prosthesis and reduce amputation pain.

\item[\textbf{TPA/TPA}] Abbreviation for tissue plasminogen activator.

\item[\textbf{TSD}] TSD is a context-dependent abbreviation that may mean total sleep deprivation, traumatic stress disorder, or another clinical term; the full term should be specified.

\item[\textbf{TSE}] TSE is a context-dependent abbreviation, including transmissible spongiform encephalopathy, the prion-disease group that includes Creutzfeldt-Jakob disease.

\item[\textbf{UL}] UL is a context-dependent abbreviation that may mean upper limit, ulnar or ulna-related, or another clinical term; the full meaning should be specified.

\item[\textbf{URG}] URG is a context-dependent abbreviation that may refer to urinary urgency or an urgent clinical priority; documentation should state the full term.

\item[\textbf{URTI}] Abbreviation for upper respiratory tract infection, commonly involving the nose, sinuses, pharynx, or larynx and often caused by viruses.

\item[\textbf{VD}] Abbreviation for venereal disease, an older term for a sexually transmitted infection.

\item[\textbf{WBS}] Abbreviation for Williams--Beuren syndrome, a genetic condition associated with characteristic facial features, developmental differences, and cardiovascular disease. In other clinical contexts, WBS may have a different meaning, so the full term should be specified.
\end{description}
